
% PIGEAN VERSION 11182025
% Template Version: Version 3 October 2023
% See section 11 of the User Manual for version history
%
%%%%%%%%%%%%%%%%%%%%%%%%%%%%%%%%%%%%%%%%%%%%%%%%%%%%%%%%%%%%%%%%%%%%%%
%%                                                                 %%
%% Please do not use \input{...} to include other tex files.       %%
%% Submit your LaTeX manuscript as one .tex document.              %%
%%                                                                 %%
%% All additional figures and files should be attached             %%
%% separately and not embedded in the \TeX\ document itself.       %%
%%                                                                 %%
%%%%%%%%%%%%%%%%%%%%%%%%%%%%%%%%%%%%%%%%%%%%%%%%%%%%%%%%%%%%%%%%%%%%%

%%\documentclass[referee,sn-basic]{sn-jnl}% referee option is meant for double line spacing

%%=======================================================%%
%% to print line numbers in the margin use lineno option %%
%%=======================================================%%

%%\documentclass[lineno,sn-basic]{sn-jnl}% Basic Springer Nature Reference Style/Chemistry Reference Style

%%======================================================%%
%% to compile with pdflatex/xelatex use pdflatex option %%
%%======================================================%%

%%\documentclass[pdflatex,sn-basic]{sn-jnl}% Basic Springer Nature Reference Style/Chemistry Reference Style


%%Note: the following reference styles support Namedate and Numbered referencing. By default the style follows the most common style. To switch between the options you can add or remove “Numbered” in the optional parenthesis. 
%%The option is available for: sn-basic.bst, sn-vancouver.bst, sn-chicago.bst%  
 
\documentclass[sn-nature,lineno]{sn-jnl}% Style for submissions to Nature Portfolio journals
%%\documentclass[sn-basic]{sn-jnl}% Basic Springer Nature Reference Style/Chemistry Reference Style
%%\documentclass[sn-mathphys-num]{sn-jnl}% Math and Physical Sciences Numbered Reference Style 
%%\documentclass[sn-mathphys-ay]{sn-jnl}% Math and Physical Sciences Author Year Reference Style
%%\documentclass[sn-aps]{sn-jnl}% American Physical Society (APS) Reference Style
%%\documentclass[sn-vancouver,Numbered]{sn-jnl}% Vancouver Reference Style
%%\documentclass[sn-apa]{sn-jnl}% APA Reference Style 
%%\documentclass[sn-chicago]{sn-jnl}% Chicago-based Humanities Reference Style

%%%% Standard Packages
%%<additional latex packages if required can be included here>

% Changing geoemtry to look better (imo)
\geometry{a4paper, margin=1in}

% Change font to arial 
\usepackage{helvet}
\renewcommand{\familydefault}{\sfdefault}

\usepackage{graphicx}%
\usepackage{multirow}%
\usepackage{amsmath,amssymb,amsfonts}%
\usepackage{amsthm}%
\usepackage{mathrsfs}%
\usepackage[title]{appendix}%
\usepackage{xcolor}%
\usepackage{textcomp}%
\usepackage{manyfoot}%
\usepackage{booktabs}%
\usepackage{algorithm}%
\usepackage{algorithmicx}%
\usepackage{algpseudocode}%
\usepackage{listings}%
\usepackage{subcaption}
\usepackage{comment}
\usepackage{tikz}
\usetikzlibrary{positioning,arrows}
\usepackage{luacode}
\usepackage{pdfpages}
\usepackage{hyperref}
%%%%

%%%%%=============================================================================%%%%
%%%%  Remarks: This template is provided to aid authors with the preparation
%%%%  of original research articles intended for submission to journals published 
%%%%  by Springer Nature. The guidance has been prepared in partnership with 
%%%%  production teams to conform to Springer Nature technical requirements. 
%%%%  Editorial and presentation requirements differ among journal portfolios and 
%%%%  research disciplines. You may find sections in this template are irrelevant 
%%%%  to your work and are empowered to omit any such section if allowed by the 
%%%%  journal you intend to submit to. The submission guidelines and policies 
%%%%  of the journal take precedence. A detailed User Manual is available in the 
%%%%  template package for technical guidance.
%%%%%=============================================================================%%%%

%% as per the requirement new theorem styles can be included as shown below
\theoremstyle{thmstyleone}%
\newtheorem{theorem}{Theorem}%  meant for continuous numbers
%%\newtheorem{theorem}{Theorem}[section]% meant for sectionwise numbers
%% optional argument [theorem] produces theorem numbering sequence instead of independent numbers for Proposition
\newtheorem{proposition}[theorem]{Proposition}% 
%%\newtheorem{proposition}{Proposition}% to get separate numbers for theorem and proposition etc.

\theoremstyle{thmstyletwo}%
\newtheorem{example}{Example}%
\newtheorem{remark}{Remark}%

\theoremstyle{thmstylethree}%
\newtheorem{definition}{Definition}%

%% New commands
% Probablity Dist function
\newcommand{\prob}{\text{Pr}}
% Variable lookup function
\newcommand{\getvar}[1]{\directlua{lookup("#1")}}

\raggedbottom
%%\unnumbered% uncomment this for unnumbered level heads

\begin{document}

\title[Article Title]{Beyond Association: Distinguishing Disease-Relevant from Disease-Associated Genes Reveals Discovery Saturation Across the Phenome}

%Dual embeddings of 6,915 traits and 18,125 genes learned from direct and indirect human genetic support
%A dense statistical map of direct and indirect genetic support between 18,125 genes and 6,915 rare and common human traits


%%=============================================================%%
%% GivenName	-> \fnm{Joergen W.}
%% Particle	-> \spfx{van der} -> surname prefix
%% FamilyName	-> \sur{Ploeg}
%% Suffix	-> \sfx{IV}
%% \author*[1,2]{\fnm{Joergen W.} \spfx{van der} \sur{Ploeg} 
%%  \sfx{IV}}\email{iauthor@gmail.com}
%%=============================================================%%

\author*[1,2]{\fnm{First} \sur{Author}}\email{iauthor@gmail.com}

\author[2,3]{\fnm{Second} \sur{Author}}\email{iiauthor@gmail.com}
\equalcont{These authors contributed equally to this work.}

\author[1,2]{\fnm{Third} \sur{Author}}\email{iiiauthor@gmail.com}
\equalcont{These authors contributed equally to this work.}

\affil*[1]{\orgdiv{Department}, \orgname{Organization}, \orgaddress{\street{Street}, \city{City}, \postcode{100190}, \state{State}, \country{Country}}}

\affil[2]{\orgdiv{Department}, \orgname{Organization}, \orgaddress{\street{Street}, \city{City}, \postcode{10587}, \state{State}, \country{Country}}}

\affil[3]{\orgdiv{Department}, \orgname{Organization}, \orgaddress{\street{Street}, \city{City}, \postcode{610101}, \state{State}, \country{Country}}}

%%==================================%%
%% Sample for unstructured abstract %%
%%==================================%%

\abstract{
Increasing GWAS power continues to yield new genetic associations, but for how many traits has gene discovery effectively saturated? For traits still yielding discoveries, how much further can GWAS alone take us and could integrating biological knowledge extend discovery beyond what genetic association can achieve? Answering these questions is critical for allocating research resources across the phenome. Here we address questions by distinguishing two concepts often conflated: \textit{disease-associated genes}, which carry variants reaching statistical significance in the population, and \textit{disease-relevant genes}, which causally influence disease when perturbed but may lack detectable genetic associations. We introduce PIGEAN, a Bayesian framework that separately quantifies direct genetic support (probability of association) and indirect genetic support (probability of relevance via shared functional annotations with associated genes). By tracing how these measures accumulate with GWAS power across {{NUM_TRAITS_SATURATION_ANALYSIS}} common phenotypes, we find that {{PERCENT_TRAITS_SATURATED}}\% show evidence of saturation for direct support, indicating diminishing returns from larger sample sizes. Among non-saturated traits, some are power-limited (more GWAS justified) while others are annotation-limited (better functional data needed to unlock indirect support). Beyond resource allocation, this decomposition also improves gene prioritization on a variety of axes: indirect support predicts future genetic associations, identifies more drug targets at comparable clinical success rates, and enriches for causal genes in rare disease patients.
}

%NOTES FOR ABSTRACT
%Novel inferences of course
%But beyond that, the presence of known and novel all in one place enables analyses (RS, gene function queries, gene lists for disease, gene lists for collection of phenotypes ) that could not be done previously
%In particular, shows shared basis of rare and common diseases, which increases further when considered at pathway level and not just genes


%%% Old abstract
% \abstract{Human genetic support increases the likelihood of therapeutic success and helps prioritize genes for further biological study. Here, we introduce the idea of indirect genetic support, which we distinguish from the traditional concept of direct support in that genetic associations for other genes support a target gene through shared gene annotations. We derive a mathemtical model, PIGEAN, that probabilistically estimates indirect genetic support from GWAS summary statisics, exome gene-level association statistics, and/or rare disease gene lists; we show that this model generalizes the idea of "bottom-up" or "convergence" based methods for gene prioritization and enables its combination with estimates of direct genetic support. We apply this model across \getvar{NUM_COMPLEX_TRAITS} complex traits, XXX rare diseases, and XXX gene annotations, assigning X\% of genes support for at least one trait. [Result from analysis of gene annotations]. [Result from analysis of how many GWAS loci have known mechanisms]. [Result from joint vs. marginal estimates and mechanisms]. [Result from drug target analysis]. These estimates of indirect support significantly expand the number of genes with links to human traits and enable probabilistic reasoning to replace ad hoc gene prioritization approaches that combine genetic associations with prior knowledge about gene function.
% }

%%================================%%
%% Sample for structured abstract %%
%%================================%%

% \abstract{\textbf{Purpose:} The abstract serves both as a general introduction to the topic and as a brief, non-technical summary of the main results and their implications. The abstract must not include subheadings (unless expressly permitted in the journal's Instructions to Authors), equations or citations. As a guide the abstract should not exceed 200 words. Most journals do not set a hard limit however authors are advised to check the author instructions for the journal they are submitting to.
% 
% \textbf{Methods:} The abstract serves both as a general introduction to the topic and as a brief, non-technical summary of the main results and their implications. The abstract must not include subheadings (unless expressly permitted in the journal's Instructions to Authors), equations or citations. As a guide the abstract should not exceed 200 words. Most journals do not set a hard limit however authors are advised to check the author instructions for the journal they are submitting to.
% 
% \textbf{Results:} The abstract serves both as a general introduction to the topic and as a brief, non-technical summary of the main results and their implications. The abstract must not include subheadings (unless expressly permitted in the journal's Instructions to Authors), equations or citations. As a guide the abstract should not exceed 200 words. Most journals do not set a hard limit however authors are advised to check the author instructions for the journal they are submitting to.
% 
% \textbf{Conclusion:} The abstract serves both as a general introduction to the topic and as a brief, non-technical summary of the main results and their implications. The abstract must not include subheadings (unless expressly permitted in the journal's Instructions to Authors), equations or citations. As a guide the abstract should not exceed 200 words. Most journals do not set a hard limit however authors are advised to check the author instructions for the journal they are submitting to.}

\keywords{keyword1, Keyword2, Keyword3, Keyword4}

%%\pacs[JEL Classification]{D8, H51}

%%\pacs[MSC Classification]{35A01, 65L10, 65L12, 65L20, 65L70}

\maketitle

\newpage
\tableofcontents
\newpage

\section{Main}\label{sec:introduction}

% TODO: Figure 1 should illustrate: (a) disease-associated vs disease-relevant distinction,
% (b) the discovery cycle concept showing saturation patterns, (c) PIGEAN framework overview

% P1: The fundamental distinction - associated vs relevant
Drug targets with human genetic support are substantially more likely to succeed in clinical trials \cite{plenge_validating_2013,minikel_refining_2024}, yet fewer than 5\% of active drug programs leverage germline genetic data \cite{minikel_refining_2024}. This gap persists in part because the field conflates two fundamentally distinct concepts: \textit{disease-associated genes}, which carry genetic variants that reach statistical significance in population studies, and \textit{disease-relevant genes}, which causally influence disease when perturbed but may lack detectable genetic associations. Disease-associated genes are a subset of disease-relevant genes (Figure~\ref{fig:main}a), constrained by statistical power \cite{wang_statistical_2019}, the presence and frequency of functional variants \cite{gazal_linkage_2017}, and the difficulty of mapping noncoding signals to effector genes \cite{schipper_demystifying_2022,costanzo_realizing_2025}. The full set of disease-relevant genes---including both core mechanistic genes and peripheral network modifiers---is what matters for therapeutic development \cite{finan_druggable_2017,jhamb_pathway_2019}, yet we lack a principled framework for distinguishing relevance from association.

% P2: The discovery cycle problem - where are we for each trait?
This distinction has practical consequences for how we allocate research resources. For any given trait, GWAS discovery follows a characteristic cycle: early stages are limited by statistical power, but eventually the rate of new discoveries saturates even as sample sizes grow \cite{abdellaoui_15_2023}. At saturation, additional GWAS investment yields diminishing returns---further progress requires identifying disease-relevant genes through other means. Yet we currently have no systematic way to diagnose where a trait sits in this discovery cycle, leading to continued investment in ever-larger biobanks for traits that may have already saturated, while neglecting the development of better functional annotations that could unlock the remaining disease-relevant genes.

% P3: Non-genetic approaches and their limitations
Non-human-genetic approaches attempt to identify disease-relevant genes beyond those captured by GWAS. Experimental studies in animal \cite{seok_genomic_2013,widjaja_inhibiting_2019} and cellular \cite{ben-david_genetic_2018,canon_clinical_2019} systems establish causal gene-phenotype links but with uncertain human relevance. Observational studies using gene expression \cite{wu_integrating_2022,rodriguez_machine_2021} or epidemiological associations \cite{barter_cholCesteryl_2011,bhatt_cardiovascular_2019} provide human context but lack causal directionality. These sources populate biomedical databases widely used for drug discovery \cite{swinney_how_2011,dornbos_evaluating_2022} and can be efficiently queried with large language models \cite{toufiq_harnessing_2023,hu_evaluation_2025}. However, gene prioritization from these sources is prone to well-documented biases toward well-studied genes, incompleteness, and circularity \cite{gillis_assessing_2013,haynes_gene_2018,stoeger_large_2018,tan_caution_2025}.

% P4: Existing integration methods are heuristic
Ideally, human genetic associations would be combined with non-genetic evidence to identify disease-relevant genes more comprehensively. Expert synthesis achieves this for individual genes \cite{mountjoy_open_2022,karjalainen_genome-wide_2024}, but genome- and phenome-wide mapping requires scalable statistical methods. Such methods exist for GWAS locus-to-gene mapping, where functional annotations enriched for associations prioritize nearby genes \cite{pers_biological_2015,weeks_leveraging_2023}, and for rare disease diagnostics, where literature mining \cite{birgmeier_amelie_2020} and interaction networks \cite{smedley_next-generation_2015} prioritize patient variants. These approaches, however, rely on heuristics such as harmonic averaging \cite{buniello_open_2025} or regression on noisy training data \cite{weeks_leveraging_2023,schipper_prioritizing_2025}, lack explicit biological models, and do not distinguish association from relevance.

% P5: Our contribution - PIGEAN framework and discovery cycle diagnosis
Here, we introduce a Bayesian framework that separately quantifies \textit{direct genetic support} (the probability that a gene is disease-associated based on genetic data) and \textit{indirect genetic support} (the probability that a gene is disease-relevant based on shared functional annotations with associated genes). Indirect support is calibrated to the same scale as direct support, enabling principled combination while preserving the distinction between association and relevance. Applied across {{NUM_TOTAL_TRAITS_MOUSE_MSIGDB:,}} traits ({{NUM_COMMON_TRAITS_MOUSE_MSIGDB:,}} common, {{NUM_RARE_TRAITS_MOUSE_MSIGDB:,}} rare) and {{NUM_TOTAL_GENE_ANNOTATIONS:,}} gene annotations, we use this framework to diagnose discovery status across the phenome, finding that {{PERCENT_TRAITS_SATURATED}}\% of common diseases show evidence of GWAS saturation---where indirect support continues to identify disease-relevant genes while direct support has plateaued. We validate indirect support on three independent axes: it predicts held-out genetic associations, prioritizes drug targets with higher clinical success rates than existing resources, and enriches for causal genes in rare disease patients. This framework provides both a diagnostic for research resource allocation and a practical tool for identifying disease-relevant genes beyond the reach of GWAS alone.

% ANALYSIS IDEA: Consider adding a "headline number" for how many additional drug targets
% indirect support surfaces compared to direct-only approaches at equivalent success rates.
% This would strengthen the translational impact claim.

% TODO: PERCENT_TRAITS_SATURATED - Calculate from the ~1000 trait saturation analysis.
% Operational definition needed: e.g., traits where slope of indirect vs direct curve < 0.1
% at current sample size, or where indirect/direct ratio exceeds threshold.


\section{Results}

% Results story (updated 01/13/2026):
% 1. Framework overview - direct vs indirect support, the PIGEAN model
% 2. Validation: indirect predicts direct (LOCO, temporal, rare disease)
% 3. Discovery Cycle: X% of traits show saturation - the headline result
% 4. Context-dependent relevance: which annotations for which contexts
% 5. Drug target prioritization: RS analysis
% 6. Rare disease gene enrichment: clinical validation

\subsection{A Framework for Estimating Indirect Genetic Support}

\begin{figure}[h!]
    \centering
    {{canva:3}}
    \caption{Conceptual overview of the PIGEAN framework. (A) Disease-associated genes identified through GWAS or Mendelian databases. (B) Disease-relevant genes share functional annotations with associated genes. (C) Graphical model of PIGEAN. (D) Genes positioned along direct and indirect support axes. (E) Downstream applications including discovery cycle diagnosis.}
    \label{fig:pigean-method}
\end{figure}

We define a \textit{disease-associated gene} as one carrying genetic variants that reach statistical significance in GWAS or that has been curated in a Mendelian database. We define a \textit{disease-relevant gene} as one for which perturbing its activity can causally influence disease risk or progression. This includes both core mechanistic genes and peripheral network modifiers, and does not require observed genetic associations. Disease-associated genes are a subset of disease-relevant genes, and our goal is to quantify both.

To estimate direct genetic support, we identified independent association signals via LD-based clumping \cite{purcell_plink_2007} and assigned signals to genes using a probabilistic nearest-gene approach that accounts for gene density and distance \cite{dornbos_evaluating_2022} (\ref{sec:methods}). To estimate indirect genetic support, we developed PIGEAN (Priors Inferred from GEne ANnotations), a Bayesian framework in which functional annotations modulate the prior probability of gene relevance, with annotation effect sizes learned from genetic data (\ref{sec:methods}). Genes sharing annotations with associated genes accumulate indirect support even when they lack direct associations. Both direct and indirect support are expressed as Bayes factors on a common scale, and indirect support exceeding 1 indicates evidence for disease relevance beyond baseline.

A key feature of PIGEAN is a spike-and-slab prior on geneset effect sizes that strictly prunes redundant or trait-irrelevant annotations. This allows us to include large heterogeneous annotation libraries representing broad swaths of current biological knowledge without inflating false positives. Genetic associations serve as the anchor, and only annotations enriched among associated genes receive non-zero effects while irrelevant or redundant annotations are regularized to zero. This design also provides a lower bound on discovery saturation. Even our best-performing annotation libraries are incomplete and cannot capture all disease-relevant biology. If indirect support plateaus despite increasing GWAS power, discovery has saturated given current knowledge. The true set of disease-relevant genes may be larger, but identifying them would require annotations that do not yet exist.

We applied PIGEAN across {{NUM_COMPLEX_TRAITS:,}} complex traits and {{NUM_RARE_DISEASE_GROUPS:,}} rare disease groups using {{NUM_TOTAL_GENE_ANNOTATIONS:,}} gene annotations spanning knockout mouse phenotypes ({{NUM_MPO_TERMS:,}} MPO terms), curated pathways ({{NUM_MSIGDB_GENESETS:,}} MSigDB gene sets), literature co-occurrence ({{NUM_PFOCR_SETS:,}} PFOCR sets, {{NUM_MESH_TERMS:,}} MeSH terms), protein interactions ({{NUM_STRING_NEIGHBORS:,}} STRING neighbors), and single-cell expression ({{NUM_SINGLE_CELL_CLUSTERS:,}} cell clusters). Model selection identified a parsimonious combination of mouse phenotypes and MSigDB as the top performer (see below), producing {{NUM_GENETRAIT_REL_PP90:,}} gene-trait relationships with posterior probability above 90\% and {{NUM_GENETRAIT_REL_PP50:,}} above 50\%. Each relationship is accompanied by provenance identifying the specific SNP associations or gene records contributing to direct support and the specific genesets contributing to indirect support.


\input{sections/indirect-support/01132026/validation_genetics_01132026}

\subsection{The Discovery Cycle: {{PERCENT_TRAITS_SATURATED}}\% of Traits Show GWAS Saturation}

\begin{figure}[t!]
    \centering
    {{fig:discovery_cycle}}
    \caption{The genetic discovery cycle across the phenome. (a) Schematic illustrating the expected discovery metric and saturation detection approach. (b-e) Exemplar traits representing different discovery stages: (b) Parkinson's disease (early, power-limited---all curves linear), (c) Asthma (mid-stage, annotation-limited), (d) Type 2 diabetes (saturated, pathway-concentrated---indirect diverges from direct), (e) Height (saturated, distributed---all curves saturate together). (f) Distribution of discovery status across {{NUM_TRAITS_SATURATION_ANALYSIS}} common traits.}
    \label{fig:discovery-cycle}
\end{figure}

To understand where research resources should be allocated across the phenome, we developed a framework for diagnosing each trait's position in the genetic discovery cycle. For any given trait, GWAS discovery follows a characteristic trajectory: early stages are limited by statistical power, but eventually the rate of new gene discoveries saturates even as sample sizes grow. Distinguishing power-limited from saturated traits is critical for research prioritization---continued GWAS investment yields diminishing returns for saturated traits, whereas power-limited traits benefit from larger cohorts.

We measure gene discovery using an \textit{expected discovery} metric that quantifies how much evidence has accumulated for each gene above a baseline prior. For any evidence type (direct, indirect, or combined), the expected discovery for gene $i$ is the posterior probability of disease relevance minus the prior probability: $D_i = P_i^{\text{post}} - P_i^{\text{prior}}$, where $P_i^{\text{post}} = \pi \cdot \text{ABF}_i / (1 + \pi \cdot \text{ABF}_i)$ and $\pi = 0.05$ is the background prior (\ref{sec:methods}). The total expected discovery $\tilde{N} = \sum_i D_i$ represents the expected number of genes identified above baseline---a continuous metric that avoids arbitrary significance thresholds while remaining interpretable as a gene count. Results were robust to the choice of background prior (Supplementary Figure~\ref{fig:prior-sensitivity}).

To trace how gene discovery accumulates with statistical power, we applied a downsampling procedure to {{NUM_TRAITS_SATURATION_ANALYSIS}} common traits with meta-analyzed GWAS summary statistics. For each trait, we simulated reduced sample sizes by inflating the standard error of each SNP's effect estimate according to the downsampled fraction $f \in \{0.1, 0.2, \ldots, 1.0\}$ (\ref{sec:methods}). At each fraction, we computed expected discovery separately for direct support ($\tilde{N}_D$), indirect support ($\tilde{N}_I$), and combined support ($\tilde{N}_C$), yielding discovery trajectories as a function of effective sample size.

We classified each trait's discovery status by comparing the fit of two models to its trajectory. A linear model implies that gene discovery continues to increase proportionally with sample size---the trait remains power-limited. An exponential saturation model implies that discovery is approaching an asymptote---the trait has reached or is entering saturation. We compared models using the Bayesian Information Criterion and further quantified the degree of saturation as the proximity to the fitted asymptote (\ref{sec:methods}).

Across {{NUM_TRAITS_SATURATION_ANALYSIS}} common traits, {{PERCENT_TRAITS_SATURATED}}\% showed evidence of GWAS saturation for direct genetic support (Figure~\ref{fig:discovery-cycle}f). The remaining traits split between power-limited ({{NUM_TRAITS_POWER_LIMITED}} traits) and annotation-limited ({{NUM_TRAITS_ANNOTATION_LIMITED}} traits). Among saturated traits, we observed two distinct architectural patterns based on the relationship between direct and indirect discovery trajectories.

Parkinson's disease exemplifies a trait in the early discovery stage (Figure~\ref{fig:discovery-cycle}b). All discovery curves---direct, indirect, combined, and the number of independent loci---increase linearly with sample size, with no evidence of saturation. The parallel trajectories indicate that both disease-associated and disease-relevant genes are being discovered at comparable rates, and continued GWAS investment will yield new loci that in turn enable identification of additional relevant genes through indirect support.

Asthma represents a mid-stage trait where direct discovery is progressing but indirect support lags (Figure~\ref{fig:discovery-cycle}c). Direct support shows early signs of curvature while the number of independent loci continues to grow linearly. Critically, indirect support accumulates more slowly than direct support, suggesting that current functional annotations incompletely capture asthma-relevant biology. For such traits, investment in disease-relevant annotations---airway epithelial expression programs or immune pathway curation---may yield greater returns than additional GWAS sample size alone.

Type 2 diabetes exemplifies a saturated trait with pathway-concentrated genetic architecture (Figure~\ref{fig:discovery-cycle}d). Both direct and indirect support show clear exponential saturation, reaching their asymptotes by approximately 40\% of the full sample size. However, combined support diverges upward from direct support, and the number of independent loci continues increasing linearly even as per-gene discovery saturates. This pattern indicates that T2D biology is organized around discrete pathways: once anchored by associated genes, these pathways propagate indirect support to functionally related genes that lack direct associations. For T2D, the discovery frontier lies in understanding pathway relationships rather than finding additional GWAS loci.

Height presents a contrasting saturation pattern (Figure~\ref{fig:discovery-cycle}e). Like T2D, both direct and indirect support saturate early---by approximately 30\% of the full sample size. Unlike T2D, all curves saturate together without divergence: direct, indirect, and combined support reach similar asymptotes. This pattern reflects height's highly polygenic, distributed genetic architecture \cite{yengo_saturated_2022}: thousands of genes contribute small effects through diverse biological processes, and shared functional annotations do not strongly connect associated genes to additional relevant genes beyond those already detected.

These results provide a principled basis for allocating resources across the phenome. For the {{NUM_TRAITS_POWER_LIMITED}} power-limited traits, continued investment in larger GWAS cohorts is justified. For the {{NUM_TRAITS_ANNOTATION_LIMITED}} annotation-limited traits, resources should be directed toward generating disease-relevant functional annotations. For the {{PERCENT_TRAITS_SATURATED}}\% of saturated traits, the discovery frontier has shifted from finding new associations to mechanistic understanding, and for pathway-concentrated traits specifically, indirect support provides a route to disease-relevant genes beyond the reach of GWAS alone.

% ANALYSIS IDEA: Temporal validation - traits classified as "power-limited" in older GWAS
% should show more new discoveries in newer GWAS than traits classified as "saturated."

% FIGURE SUGGESTION: Supplementary figure showing all ~1000 traits in 2D space
% (saturation degree vs. direct/indirect divergence), colored by disease category.


\subsection{Context-Dependent Relevance: Calibrating Annotation Sources}

\begin{figure}[t!]
    \centering
    {{fig:model_selection}}
    \caption{Model selection across geneset libraries. (a) Performance of PIGEAN models defined by different geneset library combinations on common disease validation traits, measured by relative success and mechanistic alignment. Marker size indicates model complexity. (b) Same analysis on rare disease validation traits. (c) Proportion of significant genesets from each library across trait-model pairs.}
    \label{fig:model-selection}
\end{figure}

The concept of disease relevance is inherently context-dependent: a gene relevant in the context of mouse model studies may differ from one relevant to human therapeutic intervention. Our framework operationalizes this by allowing different geneset libraries to define different relevance contexts. A model trained on mouse knockout phenotypes captures relevance as defined by mammalian loss-of-function studies; a model incorporating all available annotations captures a broader notion of relevance anchored by the full scope of current biological knowledge; and a simplified model selected for performance across diverse applications provides a practical balance of accuracy and efficiency.

To evaluate how annotation sources contribute to indirect support, we constructed multiple PIGEAN models, each defined by a specific combination of geneset libraries, and assessed their performance on held-out validation sets (Figure~\ref{fig:model-selection}a-b). Libraries included knockout mouse phenotypes ({{NUM_MPO_TERMS:,}} Mammalian Phenotype Ontology terms), curated pathways ({{NUM_MSIGDB_GENESETS:,}} MSigDB gene sets), literature co-occurrence ({{NUM_PFOCR_SETS:,}} Pathway Figure OCR sets, {{NUM_MESH_TERMS:,}} MeSH terms), protein interactions ({{NUM_STRING_NEIGHBORS:,}} STRING neighbors), and single-cell expression patterns ({{NUM_SINGLE_CELL_CLUSTERS:,}} cell type clusters).

% TODO: Add specific model comparison results
% Which combinations performed best on common vs rare?
% How much does adding more libraries help vs. hurt?

Across model selection experiments, a parsimonious model combining mouse phenotypes and MSigDB gene sets (Mouse+MSigDB) emerged as the top performer for general-purpose applications (Figure~\ref{fig:model-selection}a-b). This model achieved high relative success for drug target prioritization while maintaining mechanistic interpretability, and we adopted it for phenome-wide analysis. The spike-and-slab prior inherent to PIGEAN regularizes away uninformative annotations, enabling robust performance even when including heterogeneous or partially redundant libraries.

The contribution of each library varied systematically across trait categories (Figure~\ref{fig:model-selection}c). Mouse phenotype annotations contributed disproportionately to metabolic and developmental traits, consistent with the extensive phenotyping of these systems in model organisms. MSigDB pathway annotations showed broader utility across trait categories. Literature-derived annotations (PFOCR, MeSH) contributed more to well-studied diseases but showed signatures of ascertainment bias toward historically researched genes. These patterns highlight that the ``best'' annotation source depends on both the trait of interest and the intended application.

% ANALYSIS IDEA: Show that model selection is predictive - the best model for a trait
% category on the validation set also performs best on held-out traits in that category.

% TODO: Add quantitative results for Mouse+MSigDB model selection
% - How many traits converged?
% - What was the average number of significant genesets per trait?
% - How did complexity (number of genesets) relate to performance?


\input{sections/indirect-support/01132026/validation_pharma_01132026}

\input{sections/indirect-support/01132026/validation_clinical_01132026}

\bibliography{references}% common bib file

% Commenting out for now because weird imports (need to fix)
% % Supplementary information
\newpage
\input{sections/indirect-support/01132026/supplementary_01132026}


\end{document}

% Some notes
% - pigean indirect + direct over common
% - pigean indirect over common (A2F and/or GWAS Cat)
% - combined
% - piccolo
% - OTG - All
% - magma 
% - bar or point
% - Temporal figure of indirect support regress magma2d
% - Take temporal direct support out 2c and cite in paper


% 3x2 for figure. Panel A is RS, Panel B is 1d

% We developed pigean and applied to XXXX common traits, predicted RS, and held when holding out future whatever. 
